\begin{textsample}{POS Dim 1 – 1998 – Score 17.00 – 1998/t000200.txt}  \label{ex:f1_pos_015}
You hear that this is \textbf{the} era of environment — or \textbf{biology}, or information technology...
Well, it ' s \textbf{the} era of a lot of different things that we ' re in right now.
But one thing for sure : it ' s \textbf{the} era of change.
There ' s more change going on than ever has \textbf{occurred} in \textbf{the} history of human life on earth.
And you all sort of know it, but it ' s hard to get it so that you really understand it.
And I ' ve tried to put together something that ' s a good start for this.
I ' ve tried to show in this — though \textbf{the} color doesn ' t come out — that what I ' m concerned with is \textbf{the} little 50-year time bubble that you are in.
You tend to be interested in a generation past, a generation future — your parents, your kids, things you can change over \textbf{the} next few decades — and this 50-year time bubble you kind of move along in.
And in that 50 years, if you look at \textbf{the} population curve, you find \textbf{the} population of humans on \textbf{the} earth more than doubles and we ' re up three-and-a-half times since I was born.
When you have a new baby, by \textbf{the} time that kid gets out of high school more people will be added than existed on earth when I was born.
This is unprecedented, and it ' s big.
Where it goes in \textbf{the} future is questioned.
So that ' s \textbf{the} human part.
Now, \textbf{the} human part related to animals : look at \textbf{the} left side of that.
What I call \textbf{the} human portion — humans and their livestock and pets — versus \textbf{the} natural portion — all \textbf{the} other wild animals and just — these are vertebrates and all \textbf{the} birds, etc., in \textbf{the} land and air, not in \textbf{the} water.
How does it balance?
Certainly, 10, 000 years ago, \textbf{the} civilization ' s beginning, \textbf{the} human portion was less than one tenth of one percent.
Let ' s look at it now.
You follow this curve and you see \textbf{the} whiter spot in \textbf{the} middle — that ' s your 50- year time bubble.
Humans, livestock and pets are now 97 percent of that integrated total mass on earth and all wild nature is three percent.
We have won. \textbf{The} next generation doesn ' t even have to worry about this game — it is over.
And \textbf{the} biggest problem came in \textbf{the} last 25 years : it went from 25 percent up to that 97 percent.
And this really is a sobering picture upon realizing that we, humans, are in charge of life on earth ; we ' re like \textbf{the} capricious Gods of old Greek myths, kind of playing with life — and not a great deal of wisdom injected into it.
Now, \textbf{the} third curve is information technology.
This is Moore ' s Law plotted here, which relates to density of information, but it has been pretty good for showing a lot of other things about information technology — \textbf{computers}, their use, Internet, etc.
And what ' s important is it just goes straight up through \textbf{the} top of \textbf{the} curve, and has no real limits to it.
Now try and contrast these.
This is \textbf{the} size of \textbf{the} earth going through that same — — frame.
And to make it really clear, I ' ve put all four on one graph.
There ' s no need to see \textbf{the} little detailed words on it.
That first one is humans-versus-nature ; we ' ve won, there ' s no more gain.
Human population.
And so if you ' re looking for growth industries to get into, that ' s not a good one — protecting natural creatures.
Human population is going up ; it ' s going to continue for quite a while.
Good business in obstetricians, morticians, and farming, housing, etc. — they all deal with human bodies, which require being fed, transported, housed and so on.
And \textbf{the} information technology, which connects to our brains, has no limit — now, that is a wonderful field to be in.
You ' re looking for growth opportunity?
It ' s just going up through \textbf{the} roof.
And then, \textbf{the} size of \textbf{the} Earth.
Somehow making these all compatible with \textbf{the} Earth looks like a pretty bad industry to be involved with.
So, that ' s \textbf{the} stage out of all this.
I find, for reasons I don ' t understand, I really do have a goal.
And \textbf{the} goal is that \textbf{the} world be desirable and sustainable when my kids reach my age — and I think that ' s — in other words, \textbf{the} next generation.
I think that ' s a goal that we probably all share.
I think it ' s a hopeless goal.
Technologically, it ' s achievable ; economically, it ' s achievable ; politically, it means sort of \textbf{the} habits, institutions of people — it ' s impossible. \textbf{The} institutions of \textbf{the} past with all their inertia are just irrelevant for \textbf{the} future, except they ' re there and we have to deal with them.
I spend about 15 percent of my time trying to save \textbf{the} world, \textbf{the} other 85 percent, \textbf{the} usual — and whatever else we devote ourselves to.
And in that 15 percent, \textbf{the} main focus is on human mind, thinking skills, somehow trying to unleash kids from \textbf{the} straightjacket of school, which is putting information and dogma into them, get them so they really think, ask tough questions, argue about serious subjects, don ' t believe everything that ' s in \textbf{the} book, think broadly or creative.
They can be.
Our school systems are very flawed and do not reward you for \textbf{the} things that are important in life or for \textbf{the} survival of civilization ; they reward you for a lot of learning and sopping up stuff.
We can ' t go into that today because there isn ' t time — it ' s a broad subject.
One thing for sure, in \textbf{the} future there is an essential feature — necessary, but not sufficient — which is doing more with less.
We ' ve got to be doing things with more efficiency using less energy, less material.
Your great-great grandparents got by on muscle power, and yet we all think there ' s this huge power that ' s essential for our lifestyle.
And with all \textbf{the} wonderful technology we have we can do things that are much more efficient : conserve, recycle, etc.
Let me just rush very quickly through things that we ' ve done.
Human-powered \textbf{airplane} — Gossamer Condor sort of started me in this direction in 1976 and 77, winning \textbf{the} Kremer prize in aviation history, followed by \textbf{the} Albatross.
And we began making various odd planes and creatures.
Here ' s a giant flying replica of a pterosaur that has no tail.
Trying to have it fly straight is like trying to shoot an arrow with \textbf{the} feathered end forward.
It was a tough job, and boy it made me have a lot of respect for nature.
This was \textbf{the} full size of \textbf{the} original creature.
We did things on land, in \textbf{the} air, on water — vehicles of all different kinds, usually with some electronics or electric power systems in them.
I find they ' re all \textbf{the} same, whether its land, air or water.
I ' ll be focusing on \textbf{the} air here.
This is a solar-powered \textbf{airplane} — 165 miles carrying a person from France to England as a symbol that solar power is going to be an important part of our future.
Then we did \textbf{the} solar car for General Motors — \textbf{the} Sunracer — that won \textbf{the} race in Australia.
We got a lot of people thinking about electric cars, what you could do with them.
A few years later, when we suggested to GM that now is \textbf{the} time and we could do a thing called \textbf{the} Impact, they sponsored it, and here ' s \textbf{the} Impact that we developed with them on their programs.
This is \textbf{the} demonstrator.
And they put huge effort into turning it into a commercial product.
With that preamble, let ' s show \textbf{the} first two-minute videotape, which shows a little \textbf{airplane} for surveillance and moving to a giant \textbf{airplane}.
Narrator : A tiny \textbf{airplane}, \textbf{the} AV Pointer serves for surveillance — in effect, a pair of roving eyeglasses.
A cutting-edge example of where miniaturization can lead if \textbf{the} operator is remote from \textbf{the} vehicle.
It is convenient to carry, assemble and launch by hand.
Battery-powered, it is silent and rarely noticed.
It sends high-resolution video pictures back to \textbf{the} operator.
With onboard GPS, it can navigate autonomously, and it is rugged enough to self-land without damage. \textbf{The} modern sailplane is superbly efficient.
Some can glide as flat as 60 feet forward for every foot of descent.
They are powered only by \textbf{the} energy they can extract from \textbf{the} atmosphere — an atmosphere nature stirs up by solar energy.
Humans and soaring birds have found nature to be generous in providing replenishable energy.
Sailplanes have flown over 1, 000 miles, and \textbf{the} \textbf{altitude} record is over 50, 000 feet. \textbf{The} Solar Challenger was made to serve as a symbol that photovoltaic cells can produce real power and will be part of \textbf{the} world ' s energy future.
In 1981, it flew 163 miles from Paris to England, solely on \textbf{the} power of sunbeams, and established a basis for \textbf{the} Pathfinder. \textbf{The} message from all these vehicles is that ideas and technology can be harnessed to produce remarkable gains in doing more with less — gains that can help us attain a desirable balance between technology and nature. \textbf{The} stakes are high as we speed toward a challenging future.
Buckminster Fuller said it clearly :"there are no passengers on spaceship Earth, only crew.
We, \textbf{the} crew, can and must do more with less — much less."Paul MacCready : If we could have \textbf{the} second video, \textbf{the} one-minute, put in as quickly as you can, which — this will show \textbf{the} Pathfinder \textbf{airplane} in some flights this past year in Hawaii, and will show a sequence of some of \textbf{the} beauty behind it after it had just flown to 71, 530 feet — higher than any propeller \textbf{airplane} has ever flown.
It ' s amazing : just on \textbf{the} puny power of \textbf{the} sun — by having a super lightweight plane, you ' re able to get it up there.
It ' s part of a long-term program NASA sponsored.
And we worked very closely with \textbf{the} whole thing being a team effort, and with wonderful results like that flight.
And we ' re working on a bigger plane — 220-foot span — and an intermediate-size, one with a regenerative fuel cell that can store excess energy during \textbf{the} day, feed it back at night, and stay up 65, 000 feet for months at a time.
Ray Morgan ' s voice will come in here.
There he ' s \textbf{the} project manager.
Anything they do is certainly a team effort.
He ran this program.
Here ' s... some things he showed as a celebration at \textbf{the} very end.
Ray Morgan : We ' d just ended a seven-month deployment of Hawaii.
For those who live on \textbf{the} mainland, it was tough being away from home. \textbf{The} \textbf{friendly} support, \textbf{the} quiet confidence, congenial hospitality shown by our Hawaiian and military hosts — this is starting — made \textbf{the} experience enjoyable and unforgettable.
PM : We have real-time IR scans going out through \textbf{the} Internet while \textbf{the} plane is flying.
And it ' s exploring without \textbf{polluting} \textbf{the} stratosphere.
That ' s its goal : \textbf{the} stratosphere, \textbf{the} blanket that really controls \textbf{the} radiation of \textbf{the} earth and permits life on earth to be \textbf{the} success that it is — probing that is very important.
And also we consider it as a sort of poor man ' s stationary satellite, because it can stay right overhead for months at a time, 2, 000 times closer than \textbf{the} real GFC synchronous satellite.
We couldn ' t bring one here to fly it and show you.
But now let ' s look at \textbf{the} other end.
In \textbf{the} video you saw that nine-pound or eight-pound Pointer \textbf{airplane} surveillance drone that Keenan has developed and just done a remarkable job.
Where some have servos that have gotten down to, oh, 18 or 25 grams, his weigh \textbf{one-third} of a gram.
And what he ' s going to bring out here is a surveillance drone that weighs about 2 ounces — that includes \textbf{the} video camera, \textbf{the} batteries that run it, \textbf{the} telemetry, \textbf{the} receiver and so on.
And we ' ll fly it, we hope, with \textbf{the} same success that we had last night when we did \textbf{the} practice.
So Matt Keenan, just any time you ' re — all right — ready to let her go.
But first, we ' re going to make sure that it ' s appearing on \textbf{the} screen, so you see what it sees.
You can imagine yourself being a mouse or fly inside of it, looking out of its camera.
Matt Keenan : It ' s switched on.
PM : But now we ' re trying to get \textbf{the} video.
There we go.
MK : Can you bring up \textbf{the} house lights?
PM : Yeah, \textbf{the} house lights and we ' ll see you all better and be able to fly \textbf{the} plane better.
MK : All right, we ' ll try to do a few laps around and bring it back in.
Here we go.
PM : \textbf{The} video worked right for \textbf{the} first few and I don ' t know why it — there it goes.
Oh, that was only a minute, but I think you ' d be safe to have that near \textbf{the} end of \textbf{the} flight, perhaps.
We get to do \textbf{the} classic.
All right.
If this hits you, it will not hurt you.
OK.
Thank you very much.
Thank you.
But now, as they say in infomercials, we have something much better for you, which we ' re working on : planes that are only six \textbf{inches} — 15 centimeters — in size.
And Matt ' s plane was on \textbf{the} cover of Popular Science last month, showing what this can lead to.
And in a while, something this size will have GPS and a video camera in it.
We ' ve had one of these fly nine miles through \textbf{the} air at 35 miles an hour with just a little battery in it.
But there ' s a lot of technology going.
There are just milestones along \textbf{the} way of some remarkable things.
This one doesn ' t have \textbf{the} video in it, but you get a little feel from what it can do.
OK, here we go.
MK : Sorry.
OK.
PM : If you can pass it down when you ' re done.
Yeah, I think — I lost a little orientation ; I looked up into this light.
It hit \textbf{the} building.
And \textbf{the} building was poorly placed, actually.
But you ' re beginning to see what can be done.
We ' re working on projects now — even wing-flapping things \textbf{the} size of hawk moths — DARPA contracts, working with Caltech, UCLA.
Where all this leads, I don ' t know.
Is it practical?
I don ' t know.
But like any basic research, when you ' re really forced to do things that are way beyond existing technology, you can get there with micro-technology, nanotechnology.
You can do amazing things when you realize what nature has been doing all along.
As you get to these small scales, you realize we have a lot to learn from nature — not with 747s — but when you get down to \textbf{the} nature ' s realm, nature has 200 million years of experience.
It never makes a mistake.
Because if you make a mistake, you don ' t leave any progeny.
We should have nothing but success stories from nature, for you or for birds, and we ' re learning a lot from its fascinating subjects.
In \textbf{concluding}, I want to get back to \textbf{the} big picture and I have just two final \textbf{slides} to try and put it in perspective. \textbf{The} first I ' ll just read.
At last, I put in three sentences and had it say what I wanted.
Over billions of years on a unique \textbf{sphere}, chance has painted a thin covering of life — complex, improbable, wonderful and fragile.
Suddenly, we humans — a recently arrived species, no longer subject to \textbf{the} checks and balances inherent in nature — have grown in population, technology and intelligence to a position of terrible power.
We now wield \textbf{the} paintbrush.
And that ' s serious : we ' re not very bright.
We ' re short on wisdom ; we ' re high on technology.
Where ' s it going to lead?
Well, inspired by \textbf{the} sentences, I decided to wield \textbf{the} paintbrush.
Every 25 years I do a picture.
Here ' s \textbf{the} one — tries to show that \textbf{the} world isn ' t getting any bigger.
Sort of a timeline, very non-linear scale, nature rates and trilobites and dinosaurs, and eventually we saw some humans with caves...
Birds were flying overhead, after pterosaurs.
And then we get to \textbf{the} civilization above \textbf{the} little TV set with a gun on it.
Then traffic jams, and power systems, and some dots for digital.
Where it ' s going to lead — I have no idea.
And so I just put robotic and natural \textbf{cockroaches} out there, but you can fill in whatever you want.
This is not a forecast.
This is a warning, and we have to think seriously about it.
And that time when this is happening is not 100 years or 500 years.
Things are going on this decade, next decade ; it ' s a very short time that we have to decide what we are going to do.
And if we can get some agreement on where we want \textbf{the} world to be — desirable, sustainable when your kids reach your age — I think we actually can reach it.
Now, I said this was a warning, not a forecast.
That was before — I painted this before we started in on making robotic versions of hawk moths and \textbf{cockroaches}, and now I ' m beginning to wonder seriously — was this more of a forecast than I wanted?
I personally think \textbf{the} surviving intelligent life form on earth is not going to be carbon-based ; it ' s going to be silicon-based.
And so where it all goes, I don ' t know. \textbf{The} one final bit of sparkle we ' ll put in at \textbf{the} very end here is an utterly impractical flight vehicle, which is a little ornithopter wing-flapping \textbf{device} that — rubber-band powered — that we ' ll show you.
MK : 32 gram.
Sorry, one gram.
PM : Last night we gave it a few too many turns and it tried to bash \textbf{the} roof out also.
It ' s about a gram. \textbf{The} tube there ' s hollow, about paper-thin.
And if this lands on you, I assure you it will not hurt you.
But if you reach out to grab it or hold it, you will destroy it.
So, be gentle, just act like a wooden Indian or something.
And when it comes down — and we ' ll see how it goes.
We consider this to be sort of \textbf{the} spirit of TED.
And you wonder, is it practical?
And it turns out if I had not been — Unfortunately, we have some light bulb changes.
We can probably get it down, but it ' s possible it ' s gone up to a greater destiny up there — — than it ever had.
And I wanted to make — just — But I want to make just two points.
One is, you think it ' s frivolous ; there ' s nothing to it.
And yet if I had not been making ornithopters like that, a little bit \textbf{cruder}, in 1939 — a long, long time ago — there wouldn ' t have been a Gossamer Condor, there wouldn ' t have been an Albatross, a Solar Challenger, there wouldn ' t be an Impact car, there wouldn ' t be a mandate on zero-emission vehicles in California.
A lot of these things — or similar — would have happened some time, probably a decade later.
I didn ' t realize at \textbf{the} time I was doing inquiry-based, hands-on things with teams, like they ' re trying to get in education systems.
So I think that, as a symbol, it ' s important.
And I believe that also is important.
You can think of it as a sort of a symbol for \textbf{learning} and TED that somehow gets you thinking of technology and nature, and puts it all together in things that are — that make this conference, I think, more important than any that ' s taken place in this country in this decade.
Thank you.

% matched lemmas: airplane, altitude, biology, cockroach, computer, conclude, crude, device, friendly, inch, learning, occur, one-third, pollute, slide, sphere, the
\end{textsample}
